\lektion{Messung, Eindeutigkeit und Sensorort}{befund123}{Nachweise}{Rückblick · Kapitel 1 bis 3}
\subt{1 · Wie genau lässt sich eine Geschwindigkeit aus Wegdaten bestimmen?}
Zwei Messungen im Abstand \(h\) haben jeweils höchstens \(\varepsilon\) Ortsfehler. Die Beschleunigung sei dem Betrag nach durch \(M\) begrenzt. Der rückwärts gebildete Differenzenquotient weicht von der momentanen Geschwindigkeit um höchstens \(Mh/2+2\varepsilon/h\) ab. Für \(M,\varepsilon>0\) wird diese Grenze bei \(h=2\sqrt{\varepsilon/M}\) am kleinsten. Ohne Beschleunigungsgrenze sind beliebig verschiedene Endgeschwindigkeiten mit denselben zwei Wegmessungen vereinbar: Addieren Sie zur Bahn ein beliebiges Vielfaches von \(t(t-h)\). Zur Bewertung eines realen Sensors fehlen gemessene Werte für \(M\) und \(\varepsilon\).
\subt{1 · Gilt die Pendelgrenze auch mit Dämpfung oder Stift?}
Schon die linearisierte Gleichung \(\varphi''+2\gamma\varphi'+\omega_0^2\varphi=0\) besitzt für \(\gamma<\omega_0\) die Periode \(2\pi/\sqrt{\omega_0^2-\gamma^2}\). Bei \(\gamma\ge\omega_0\) entfällt die periodische Schwingung. Ein Stift verkürzt während des Kontakts die wirksame Fadenlänge. Die Einprozentgrenze gilt für das freie, ungedämpfte Modell.
\subt{2 · Alle klassischen Lösungen der singulären Gleichung}
Für \(q\in C^2(I)\) auf einem Intervall gilt \(q(q''+1)=0\). Jede Lösung ist entweder identisch null oder auf ganz \(I\) von der Form \(q(t)=-t^2/2+At+B\).
\proofpart{Wo \(q\ne0\), gilt \(q''=-1\), und durch Stetigkeit gilt das auch im Abschluss dieser Menge. Im Inneren der Nullmenge gilt \(q''=0\). Damit nimmt das stetige \(q''\) höchstens die Werte \(-1\) und 0 an. Auf einem Intervall ist es konstant. Bei \(q''=0\) erzwingt die Gleichung \(q=0\), sonst ergibt zweifache Integration die Parabel. Beide Formen erfüllen die Gleichung.}
Bei \(q(0)=q'(0)=0\) sind sowohl die Nulllösung als auch \(-t^2/2\) zulässig. Andere Anfangspaare bestimmen genau eine Parabel.
\subt{3 · Welcher Sensor sieht die zweite thermische Mode?}
Für zwei symmetrische Speicher seien Mittelwert \(a\) und Differenz \(b\) die Moden, also \(T_1=a+b,\ T_2=a-b\). Ein Sensor mit \(y=sT_1+(1-s)T_2\) misst \(a+(2s-1)b\). In der Mitte verschwindet die Differenz vollständig. Unter gleicher Messunsicherheit und bekanntem Mittelwert sind die Endpunkte \(s=0,1\) am empfindlichsten für \(b\). Der beste Ort an einem konkreten Bauteil verlangt dessen Leitwerte und Messrauschen.
\srcs{\see{schluss1}{Messmodell} · \see{schluss2}{Implizite Gleichungen} · \see{schluss3}{Systemtheorie}.}

\lektion{Gilt die Fehlerschranke auch zwischen Zeitpunkten?}{befund4}{Nachweise}{Rückblick · Kapitel 4}
Der steife Folgeregler hat die exakte Lösung \(y(t)=\cos t-e^{-1000t}\). Die dichte Ausgabe eines RK45-, Radau- oder BDF-Schritts ist hier ein Polynom \(P(t)\) vom Grad höchstens fünf. Wir können seinen größten Fehler auf dem \emph{ganzen} Schritt nach oben begrenzen. Dafür brauchen wir die vom Solver gespeicherten Polynomkoeffizienten.
\subt{Von sechs Prüfstellen zur ganzen Kurve}
Teile einen Schritt in Intervalle der Breite \(h\). Interpoliere auf jedem Teilintervall den Fehler \(e=P-y\) an sechs äquidistanten Punkten durch ein Polynom \(B\) vom Grad fünf. Seine Bernstein-Darstellung ist eine gewichtete Summe von sechs Koeffizienten, deren Basisgewichte überall nichtnegativ sind und zusammen eins ergeben. Also gilt \(|B|\le\max_j |b_j|\). Die Interpolationsrestformel ergibt ab dem linken Rand \(a\)
\[
 \sup |P-y|\le \max_j|b_j|+
             \frac{h^6}{720}\bigl(1+1000^6e^{-1000a}\bigr).
\]
Denn \(P^{(6)}=0\); das Produkt der sechs Abstände zu den Stützstellen ist höchstens \(h^6\). Der Rechenlauf wertet auch die Koeffizienten und den Restterm mit nach außen gerundeter Intervallarithmetik aus.
{\tblstart\begin{tabularx}{\linewidth}{@{}Yrr@{}}\toprule
\hd{Solver}&\hd{Obere Grenze bei \(\mathrm{rtol}=10^{-3}\)}&\hd{bei \(\mathrm{rtol}=10^{-6}\)}\\\midrule
RK45&\(2{,}22\cdot10^{-3}\)&\(1{,}19\cdot10^{-6}\)\\
Radau&\(2{,}04\cdot10^{-3}\)&\(2{,}30\cdot10^{-5}\)\\
BDF&\(5{,}44\cdot10^{-4}\)&\(1{,}10\cdot10^{-6}\)\\\bottomrule\end{tabularx}\par}
Die groben Läufe überschreiten das Ziel \(10^{-4}\) bereits an einzelnen Prüfstellen. Für die drei feinen Läufe liegen die oberen Schranken auf dem gesamten Intervall von 0 bis 2 s unter dem Ziel. Die Aussage bezieht sich auf die reellen Polynome aus den gespeicherten Binärkoeffizienten von SciPy 1.17.0. Die erneute Gleitkommaauswertung an einem Gerät kann zusätzlich runden.
\notebox{Übertragen}{Eine garantierte obere Grenze oberhalb des Ziels allein beweist keinen Fehlschlag. Welcher Befund genügt? Ein einziger nachweislich zu großer Punktfehler. Die Datei enthält zu jedem Lauf beide Grenzen.}
\srcs{\see{tracking}{Ursprünglicher Vergleich} · \texttt{kapitel.json}, Abschnitt \texttt{dense}.}

\lektion{C21 · Bleibt die Diodenspannung vor der Falte?}{befunddiode}{Beweis}{Rückblick · Kapitel 5 und 10}
Für die Schaltung aus Kapitel 10 mit Nullstart und sinusförmiger Einvoltquelle liegt die Falte bei \(z_*=0{,}306158\unitm V\). Wir zeigen für den vom Nullstart erreichten Niedrigstromast, dass \(z\) unter \(0{,}1\unitm V\) bleibt. Damit bleibt die algebraische Gleichung auf diesem Ast invertierbar.
\subt{Ein Widerspruch beim ersten Erreichen}
Setze \(b=u_C+R_2i\) und \(d=i_D\), dann ist \(z=b+R_2d\). Mit \(X=(i,u_C)^\top\) lautet der lineare Schaltungsteil
\[
 X'=AX+B_su_0+B_dd,\quad
 A=\begin{pmatrix}-(R_1+R_2)/L&-1/L\\1/C&0\end{pmatrix},
 \quad B_s=\binom{1/L}{0},\quad B_d=\binom{-R_2/L}{1/C}.
\]
Für \(c^\top=(R_2,1)\) hat jeder Übertragungskern \(c^\top e^{At}B\) die Form \(\sum_j r_j e^{\lambda_jt}\), wobei die beiden verschiedenen Eigenwerte negativ sind. Die Residuen folgen direkt aus \(r_j=c^\top(A-\lambda_\ell I)B/(\lambda_j-\lambda_\ell)\), \(\ell\ne j\).
\proofpart{Für die Sinusquelle beschränkt Integration der komplexen Exponentialfunktion jeden Beitrag durch \(2|r_j|/|\mathrm i\omega-\lambda_j|\). Für beliebigen Diodenstrom liefert die Integration des Betrags \(\sum_j |r_j|/(-\lambda_j)\) als Verstärkungsgrenze. Mit den Schaltungswerten ergeben sich weniger als \(0{,}05915\unitm V\) bzw. \(981\,\Omega\). Bis zum ersten möglichen Erreichen von \(0{,}1\unitm V\) ist \(0\le d\le I_S(e^{0{,}1/U_T}-1)\). Daraus folgt \(|b|<0{,}05960\unitm V\). Zum Erreichen wäre aber \(b=0{,}1-R_2 I_S(e^{0{,}1/U_T}-1)>0{,}09999\unitm V\) nötig.}
Im leitenden Ast ist somit \(g'(z)>0{,}99963\); im Sperrast gilt \(g'=1\). Beide Äste schließen bei null stetig an. Damit entfällt für diese Parameter und diesen Start die bisher fehlende Fortsetzungsgarantie. Eine unabhängige Rechnung mit der reellen Hauptverzweigung der Lambert-\(W\)-Funktion fand in \(500\,\mu\mathrm s\) die 149 späteren Nulldurchgänge. An 6\,013 gesonderten Ausgabezeiten lag das größte \(z\) bei \(0{,}02212\unitm V\) oder darunter.
\srcs{\see{diodenast}{Kennlinie und Falte} · \texttt{kapitel.json}, Abschnitt \texttt{diode}.}

\lektion{C19R · Positivität und Fehler an der Messstelle}{befundtransport}{Nachweise}{Rückblick · Kapitel 6 und 11}
Bei nichtnegativen Zellwerten liegt der im Code rekonstruierte östliche Konzentrationswert zwischen null und dem Doppelten des Zellwerts \(c_i\). Das folgt aus der Begrenzung der Steigung durch das Doppelte des jeweiligen Nachbarsprungs. Der von Westen einströmende Wert ist ebenfalls nichtnegativ. Diffusion bringt nichtnegative Nachbarbeiträge, eine positive Quelle fügt Masse hinzu. Für einen Euler-Teilschritt genügt deshalb
\[
 h\left(2v/\Delta x+2D_x/\Delta x^2+2D_y/\Delta y^2+
 \lambda_{\max}\right)\le 1 .
\]
Hier sind \(v,D_x,D_y\) bereits um die Retardation skalierte Größen und \(\lambda_{\max}\) begrenzt die lokale Abbaurate. Der verbleibende Koeffizient von \(c_i\) ist dann nichtnegativ. Die örtliche Quellrate wird durch das Zellvolumen geteilt; Multiplikation mit Porenanteil, Dicke und Zellfläche ergibt bei \(R=1\) auf jedem getesteten Gitter \(2\unitm{mg/s}\).
\subt{Wie verhält sich das tatsächlich verwendete RK4?}
In acht Varianten des Zehntageversuchs wurden alle Zwischenstufen und Endwerte auf ihr Minimum geprüft; es lag jeweils bei null. Eine unbedingte RK4-Positivitätsgarantie ist falsch. Auf \(9\times9\) Zellen mit \(0{,}25\unitm m\) Abstand, schachbrettartigem Null-Eins-Start, Quelle außerhalb und \(h=200\,000\unitm s\) erzeugt ein einziger Schritt einen Wert unter \(-12\,535\). Die Euler-Bedingung oben ist eine bewiesene Aussage für Euler.
\subt{Was lässt sich am Brunnen garantieren?}
Bei \(x=5\unitm m,\ y=0\) nach zehn Tagen liegt die unendliche Punktquellenreferenz zwischen \(45\,014{,}84\) und \(45\,028{,}41\unitm{mg/m^3}\). Der Alterskern steigt bis zu einem Maximum und fällt dann; seine totale Variation ist höchstens zweimal die Gipfelhöhe. Damit ist der Fehler der Trapezsumme durch Schrittweite mal totale Variation eingeschlossen. Die Rechnung nutzt 65\,536 Intervalle.
Das \(0{,}25\unitm m\)-Gitter liefert \(45\,440{,}70\), das \(0{,}125\unitm m\)-Gitter im erweiterten Gebiet \(45\,085{,}52\unitm{mg/m^3}\). Die gespeicherten Brunnenwerte liegen damit garantiert innerhalb von \(0{,}95\,\%\) bzw. \(0{,}16\,\%\) der kontinuierlichen Referenz. Auf dem Viertelmetergitter bewirkt Halbierung des Zeitschritts rund \(0{,}006\unitm{mg/m^3}\) Unterschied. Eine breiter verteilte Quelle gleicher Masse verschiebt den Wert auf \(44\,212{,}64\). Am Quellpunkt selbst ist das zweidimensionale Altersintegral wegen \(1/a\) singulär; eine globale punktweise Fehlerschranke über alle Orte wäre schon deshalb ungeeignet.
\srcs{\see{c19integral}{Kern und Einheiten} · \texttt{transport.json} · \texttt{gegenproben.json}.}

\lektion{Ereignisse · Vier Fragen mit verschiedenen Antworten}{befundereignis}{Nachweise}{Rückblick · Kapitel 7 bis 10}
\subt{7 · Kann endlich häufiges Abfragen jedes Ereignis erkennen?}
Nein. Seien \(t_i\) alle besuchten Stellen eines Detektors beim konstanten Signal eins. Das glatte Signal \(1-A\prod_i(t-t_i)^2\) hat an jeder dieser Stellen denselben Wert und dieselbe erste Ableitung. Durch ein genügend großes \(A\) wird es zwischen zwei Stellen negativ. Es gibt also ununterscheidbare Abfragen mit verschiedenen Nullstellen. Eine Garantie erfordert etwa eine geeignete Monotoniegrenze, eine Einschließung auf dem ganzen Intervall oder eine auswertbare Modellstruktur.
\subt{8 · Bleiben drei C05-Schaltungen bei veränderten Schwellen?}
Ja, für die oberen Schwellen \(5{,}8\pm0{,}005\) und unteren \(2{,}5\pm0{,}005\) bei den Startwerten des Lehrversuchs. Der exakte Phasenfluss hat die Form \(E+A e^{-c_3t}+B e^{-c_1t}\). Seine Ableitung besitzt höchstens eine Nullstelle, da eine Quotientenbildung der beiden Exponentialterme eine streng monotone Funktion liefert. Start und Gleichgewicht liegen auf entgegengesetzten Seiten der aktiven Schwelle. Daher ist die erste positive Durchquerung eindeutig.
Mit nach außen gerundeten Intervallen werden beide Vorzeichen eingeklammert und der Zustandsbereich an die nächste Phase übergeben. Die vier Zeiten liegen in \([1{,}0946;1{,}1227]\), \([2{,}0961;2{,}1654]\), \([3{,}0063;3{,}1048]\), \([4{,}0078;4{,}1476]\unitm s\). Jede Durchquerung hat gerichtete Steigung größer als \(0{,}2307\) pro Sekunde. Bis \(3{,}5\unitm s\) liegen folglich genau drei Schaltungen.
\subt{9 · Wie viel Flugzeit unterschlägt eine Ruhegrenze?}
Sei \(v\le v_{\min}\) die gerade erreichte \emph{Abflug}geschwindigkeit, \(g>0\) und \(0\le e<1\) der Rückprallfaktor. Die verbleibenden Flüge dauern \(2v/g,2ev/g,\ldots\). Ihre Summe ist \(2v/[g(1-e)]\), die größte ausgelassene Höhe \(v^2/(2g)\). Ersetzt man \(v\) durch \(v_{\min}\), erhält man garantierte Obergrenzen. Eine auf die Ankunftsgeschwindigkeit bezogene Regel muss erst mit \(e\) auf den nächsten Abflug umgerechnet werden.
\subt{10 · Ist der C07-Wechsel auch kartesisch konsistent?}
Ein zweiter Integrator führt Ort und Geschwindigkeit mit der aus der Bindung abgeleiteten Beschleunigung. Beim Stiftwechsel versetzt er den Drehpunkt und übernimmt die physische Geschwindigkeit. Über zehn Sekunden wurden alle acht Wechsel geprüft. Längenfehler blieben unter \(1{,}18\cdot10^{-12}\unitm m\), \(|(r-p)\cdot v|\) unter \(5{,}33\cdot10^{-12}\unitm{m^2/s}\). Die größte Zeitabweichung von der Winkelrechnung betrug \(3{,}77\cdot10^{-11}\unitm s\). Der Impulssprung ist beim verwendeten Reset konstruktiv null.
\srcs{\texttt{kapitel.json}, Abschnitte \texttt{thresholds} und \texttt{cartesian} · \see{befunddiode}{C21-Astbeweis}.}

\lektion{Wählt der schnellste Solver das nächste Modell?}{befundwahl}{Labor}{Rückblick · Kapitel 12}
Ein Solver ist für eine neue Modellfamilie auszuwählen, bevor wir deren Zeiten kennen. Als einfache Regel betrachten wir die 17 Einstellungen, die für alle acht archivierten C-Modelle vorliegen. Wir lassen jeweils ein Modell vollständig außen vor. Auf den sieben übrigen Modellen minimieren wir zuerst die Zahl der Fehlschläge und anschließend die mittlere Laufzeit relativ zur schnellsten erfolgreichen Einstellung je Modell. Ein Gleichstand wird nach dem Methoden- und Einstellungsnamen entschieden.
\subt{Was passiert am zurückgehaltenen Modell?}
In allen acht Durchgängen wählt diese Regel RK45. Das Genauigkeitsziel erreicht er an sieben zurückgehaltenen Modellen; an C05 scheitert er. In C21 dauert sein erfolgreicher Lauf rund \(10{,}29\)-mal so lange wie der schnellste zugelassene. Das Verfahren benutzt die bereits bekannte Studie rückblickend, hält jedoch die Testzeilen im jeweiligen Auswahlalgorithmus zurück.
\subt{Zwei zusätzliche Systeme mit bekannten Lösungen}
Als separate Gegenprobe fixieren wir dieselbe Wahl und rechnen ein erzwungenes Duffing-System mit exakter Schwingung und ein steifes gekoppeltes System mit exaktem Verlauf \((\cos t,e^{-t})\). Für alle drei Solver gelten \(\mathrm{rtol}=10^{-6}\), \(\mathrm{atol}=10^{-8}\), maximal \(10^{-4}\) gemessener Fehler auf 1\,001 Ausgabezeiten und höchstens 8\,000 Aufrufe der rechten Seite im Laufprotokoll.
{\tblstart\begin{tabularx}{\linewidth}{@{}Yrrr@{}}\toprule
\hd{Modell}&\hd{RK45}&\hd{Radau}&\hd{BDF}\\\midrule
Erzwungene Schwingung&98&208&90\\
Schnelle gekoppelte Mode&50\,300&148&89\\\bottomrule\end{tabularx}\par}
Zellen nennen \texttt{nfev}; Faktorisierungen stehen separat in der Datei. RK45 überschreitet beim gekoppelten System das Aufrufbudget. Radau und BDF erfüllen hier beide Ziele. Für eine allgemeingültige Solverwahl reichen frühere Siege nicht. Eine nutzbare Regel muss ihre Modellfamilie, die Fehlergröße und das Budget mit angeben.
\notebox{Welche Frage verlangt Außendaten?}{Für die beiden realen Sensorfragen aus Kapitel 1 und 3 gibt es analytische Bedingungen und Gegenbeispiele. Eine reale Genauigkeit oder ein bester Montageort lassen sich erst mit kalibrierter Messkette, Störpegel und Daten vom Bauteil bestimmen. Das Rechenlabor enthält ausschließlich simulierte Daten.}
\srcs{\repo{docs/evidence/qss-selection.csv}{Archiv mit 520 Konfigurationen} · \texttt{kapitel.json} · \texttt{gegenproben.json}.}

\lektion{C07 · Zwei Bindungen in jedem Rechenschritt}{befundrattle}{Labor}{Rückblick · Kapitel 10}
Ein Pendel der Länge \(\ell\) mit Drehpunkt \(p\) muss zwei Bedingungen erfüllen. Mit \(s=r-p\) lauten sie \(s\cdot s=\ell^2\) und \(s\cdot v=0\). Die RATTLE-Methode bestimmt die Zwangskorrekturen bei Ort und Geschwindigkeit getrennt. Für Dämpfung verwenden wir je einen exakten halben Dämpfungsschritt vor und nach dem konservativen Teilschritt.
\subt{Wie lässt sich die Ortsbindung lösen?}
Nach halber Dämpfung, halbem Schwerkraftstoß und freier Drift ergäbe sich ein Versuchspunkt \(w\) relativ zum Drehpunkt. Der Zwangsstoß wirkt in der bisherigen Seilrichtung \(n=s/\ell\). Der orthogonale Anteil \(w_\perp=w-(w\cdot n)n\) bleibt beim Korrigieren unverändert. Die kontinuierlich anschließende Lösung des quadratischen Ortszwangs ist
\[
 s^+=w_\perp+\sqrt{\ell^2-\|w_\perp\|^2}\,n .
\]
Anschließend bestimmen die korrigierte Ortsdifferenz, der zweite halbe Schwerkraftstoß und eine orthogonale Projektion auf die neue Tangente die Geschwindigkeit. Beim Erreichen der Stiftrichtung klammert eine Teilung des Schritts die Zeit ein. Der neue Drehpunkt muss auf derselben Strahlgeraden liegen, damit die aktuelle Position und Tangente zu beiden Kreisen passen.
{\tblstart\begin{tabularx}{\linewidth}{@{}Yrr@{}}\toprule
\hd{Schrittweite}&\hd{erkannte Wechsel}&\hd{größter Zeitfehler}\\\midrule
\(0{,}0100\unitm s\)&8&\(474{,}4\,\mu\mathrm s\)\\
\(0{,}0050\unitm s\)&8&\(115{,}4\,\mu\mathrm s\)\\
\(0{,}0025\unitm s\)&8&\(28{,}7\,\mu\mathrm s\)\\\bottomrule\end{tabularx}\par}
Referenz ist die separat integrierte kartesische Beschleunigungsform, die bereits an der Winkelrechnung geprüft wurde. In allen drei RATTLE-Läufen blieb der maximale Längenfehler unter \(4{,}45\cdot10^{-16}\unitm m\). Beim Halbieren der Schrittweite schrumpft hier der Zeitfehler um etwa den Faktor vier. Dass die Bindung bei allen drei Gittern auf Rundungsniveau erfüllt ist, macht die Unterschiede der Ereigniszeiten besonders gut sichtbar.
\srcs{\see{befundereignis}{Kontrolle der Moduswechsel} · \texttt{rattle.json} · \texttt{constrained\_pendulum.py}.}
